\chapter{Experimental Doxonomics}
\label{ch:experimental}

\epigraph{The great tragedy of science---the slaying of a beautiful hypothesis by an ugly fact.}{Thomas Henry Huxley}

\noindent
Observational methods identify correlations; experiments identify causes.
This chapter develops experimental designs for testing doxonomic hypotheses in the laboratory and the field.

\section{Cooperation Games as Belief Measurement}
\label{sec:ch19-games}

\begin{definition}[Doxonomic Game]
\label{def:doxonomic-game}
\index{doxonomic game}
A \emph{doxonomic game} is an experimental economic game (public goods, dictator, trust, ultimatum) administered after exposure to a doxonomic treatment (narrative about human nature, institutional framing, etc.).
The outcome variable is the degree of cooperative or self-interested behavior, interpreted as a \emph{behavioral proxy} for internalized belief.
\end{definition}

\textcite{fischbacher2001} found that most people are conditional cooperators---not the selfish maximizers predicted by standard economic models.
\textcite{henrich2001} showed this holds cross-culturally.
The doxonomic question is: does exposure to selfishness narratives \emph{reduce} cooperation?

\section{Priming and Framing Experiments}
\label{sec:ch19-priming}

\begin{model}[Narrative Priming]
\label{mod:priming}
\index{priming experiment}
Randomly assign subjects to:
\begin{itemize}[nosep]
  \item \textbf{Treatment}: read a text asserting that human nature is fundamentally self-interested,
  \item \textbf{Control}: read a neutral text or one asserting cooperative human nature.
\end{itemize}
Then measure behavior in a public goods game.
The treatment effect $\hat{\tau} = \bar{Y}_{\text{treat}} - \bar{Y}_{\text{control}}$ estimates the causal effect of a single narrative exposure on cooperative behavior.
\end{model}

\section{Norm Perception Studies}
\label{sec:ch19-norms}

\begin{definition}[Perceived Social Norm]
\label{def:perceived-norm}
\index{perceived norm}
The \emph{perceived social norm} for behavior $x$ is agent $j$'s belief about the population's behavior: $\hat{\mu}_j(x)$.
Doxonomic systems can shift behavior by shifting norm perceptions---making people think others are more selfish than they actually are.
\end{definition}

\begin{proposition}[Norm Misperception Effect]
\label{prop:norm-misperception}
If agents are conditional cooperators (cooperate when they believe others cooperate), then a doxonomic intervention that increases $\hat{\mu}_j(\text{selfish})$ without changing actual behavior will reduce cooperation in subsequent interactions.
The belief becomes self-fulfilling.
\end{proposition}

\section{Ethical Limits}
\label{sec:ch19-ethics}

Experimental doxonomics faces distinctive ethical constraints:
\begin{itemize}[nosep]
  \item exposing subjects to false narratives about human nature may have lasting effects,
  \item deception about the social environment may erode trust,
  \item experiments that increase selfishness ``work'' by making people worse off.
\end{itemize}

These constraints limit the scope of laboratory doxonomics but do not eliminate it.
Natural experiments and field studies with minimal deception remain viable.

\section{Notes and References}
\label{sec:ch19-notes}

Public goods experiments are reviewed in \textcite{fischbacher2001}.
Cross-cultural behavioral games are in \textcite{henrich2001}.
Priming effects on economic behavior are studied in the behavioral economics literature; see \textcite{kahneman2011} for an overview.
On norm misperception, see \textcite{cialdini2006}.
Experimental ethics in social science is discussed in \textcite{salganik2018}.

\section{Exercises}

\begin{exercise}
Design a priming experiment to test whether exposure to ``greed is good'' narratives increases defection in a public goods game.
Specify sample size, treatment arms, and outcome measures.
\end{exercise}

\begin{exercise}
A norm perception study finds that 80\% of people believe ``most people are selfish,'' but behavioral games show only 30\% defect.
Compute the norm misperception gap and predict its doxonomic consequences.
\end{exercise}

\begin{exercise}
Discuss the ethics of an experiment that exposes children to competing anthropological narratives (selfish vs.\ cooperative) and measures long-term behavioral effects.
Under what conditions (if any) would this be justifiable?
\end{exercise}
